\chapter{Related Works}

Developing a robust Advanced Driver Assistance System (ADAS) requires a multidisciplinary approach, bridging the gap between modern artificial intelligence and automotive-grade software engineering. To establish a solid foundation for the proposed red-light warning prototype, a comprehensive review of existing literature and technological frameworks is essential. The following sections explore key advancements across four primary domains that directly influence the system's architecture. Initially, the discussion examines current methodologies in traffic light detection, highlighting the transition from traditional image processing to deep learning-based recognition. Subsequently, the practical challenges of deploying these computationally demanding perception models onto hardware-constrained edge devices are analyzed. The focus then shifts to the AUTOSAR Adaptive Platform, which provides the necessary service-oriented infrastructure for modern automotive software. Finally, the integration of GNSS-based spatial positioning is reviewed as a strategic mechanism to trigger and support visual perception pipelines. Together, these areas establish the necessary context to justify the technological decisions made throughout the project's implementation.

\section{Traffic Light Detection}

Traffic light detection is a fundamental perception task for Advanced Driver Assistance Systems (ADAS) and autonomous driving because traffic signal interpretation directly affects vehicle behavior at intersections. Traditional image-processing approaches rely on handcrafted features such as color thresholding, shape analysis, and region filtering. However, these approaches are sensitive to illumination changes, occlusion, background clutter, and camera viewpoint. Recent studies therefore increasingly adopt deep learning-based object detection to localize traffic lights and classify their signal states.

Nguyen et al.~\cite{Nguyen2025TrafficLightMonitoring} present a traffic light monitoring system using artificial intelligence and show that deep learning-based recognition can improve robustness in practical traffic environments. Their work is relevant because it confirms the suitability of AI-based perception for traffic light monitoring. However, it mainly focuses on detection and monitoring performance, while the integration of perception results into an automotive software architecture and an embedded ADAS warning pipeline is less emphasized.

Lee et al.~\cite{Lee2022TrafficLightRecognition} propose a traffic light recognition system for autonomous driving vehicles using a monocular camera combined with Intelligent Transportation System (ITS) information. This work is closely related to the motivation of this project because it shows that camera-only perception may not be sufficient when the traffic light is far away, visually small, or difficult to detect early enough. By incorporating external contextual information, the system can support recognition before visual perception becomes fully reliable. This idea motivates the GPS-assisted triggering strategy in this project, where the vehicle position and traffic-light node database are used to determine when the AI camera pipeline should be activated.

Countdown timer recognition is another related perception problem. Goodfellow et al.~\cite{goodfellow2014multidigitnumberrecognitionstreet} demonstrate that deep convolutional neural networks can recognize multi-digit numbers from real-world street-view imagery. Although their work is not specific to traffic light countdown timers, it provides a strong reference for digit recognition in unconstrained outdoor scenes. This supports the timer OCR component of the proposed system, where countdown values associated with detected traffic lights are recognized and displayed to the user.

\section{Deploying Model on Hardware-Constrained Devices}

A major gap between research prototypes and practical ADAS systems is the deployment of perception models on embedded hardware. Object detection models are often evaluated on desktop GPUs, but in-vehicle systems must satisfy stricter constraints in terms of CPU usage, memory consumption, power consumption, and latency. For a driver warning function, average accuracy alone is insufficient; the model must also provide stable frame rate and bounded processing delay.

MobileNetV4~\cite{qin2024mobilenetv4universalmodels} is engineered specifically for resource-constrained ecosystems, offering an optimal balance between accuracy and computational efficiency. During the first phase of this project, we leveraged this architecture to develop a highly optimized YOLO11n-MobileNetV4 model---formally documented in our IEEE SAS 2026 submission (see \hyperref[chap:append]{Appendix}). Building directly upon the theoretical foundations and structural optimizations established in that paper, the current prototype advances the design by proposing a novel {YOLO26-MobileNetV4} model because YOLO26 is better than YOLOv11 \cite{yolo26}. This evolution aims to further minimize inference latency and maximize real-time feasibility on edge devices.

Goodfellow et al.~\cite{goodfellow2014multidigitnumberrecognitionstreet} also provide useful background for embedded perception because their CNN-based digit recognition approach demonstrates how a focused neural network can solve a specific recognition subtask. In this project, the traffic light detector and the countdown timer recognizer are separated so that the timer OCR task can be handled by a lightweight CNN model rather than by a large general-purpose detection network. This modular design supports the goal of maintaining real-time performance on constrained hardware.

\section{AUTOSAR Adaptive Platform}

The AUTOSAR Adaptive Platform is designed for high-performance automotive applications that require dynamic service-oriented communication, lifecycle management, and execution on POSIX-based operating systems. AUTOSAR documentation describes the Adaptive Platform as a complement to Classic AUTOSAR~\cite{autosar_layered_arch}, targeting applications such as ADAS, automated driving, infotainment, and high-performance computing workloads~\cite{autosar_platform_design,autosar_sw_arch}. This makes it an appropriate architectural foundation for integrating perception, communication, and display functions in the proposed prototype.

Communication Management is a central functional cluster in Adaptive AUTOSAR. It defines how Adaptive Applications communicate through service-oriented interfaces and network bindings such as SOME/IP~\cite{autosar_someip} and DDS~\cite{autosar_com_mgmt, autosar_dds}. Execution Management defines how processes are started, supervised, and coordinated using execution manifests and machine states~\cite{autosar_exec_mgmt}. Log and Trace provides standardized runtime observability for debugging and validation~\cite{autosar_log_trace}. These functional clusters are relevant to this project because the prototype separates the Provider and Client applications, uses middleware-based communication, reports execution state, and relies on logging to evaluate runtime behavior.

The AUTOSAR methodology further defines a configuration-driven workflow in which high-level architecture, service interfaces, application design, machine configuration, and deployment artefacts are refined into executable software~\cite{AUTOSAR_TR_AdaptiveMethodology_R22-11}. PopcornSAR provides AUTOSAR tooling and product support aligned with this configuration-oriented workflow~\cite{popcornsar_intro}. These references form the methodological basis for using Autosar.io to configure the prototype, generate application artefacts, and maintain consistency between design and deployment.

\section{Spatial Positioning with GNSS}

Spatial positioning provides an important complementary layer for traffic light warning systems. A camera can only detect what is visible in its field of view, and its reliability decreases when the traffic light is far away, small, occluded, or affected by lighting conditions. At higher vehicle speeds, waiting until the traffic light is clearly visible may leave insufficient time for warning generation. Therefore, a supporting positioning mechanism is needed to determine whether the vehicle is approaching a relevant traffic-light node before the camera-based detector becomes reliable.

Lee et al.~\cite{Lee2022TrafficLightRecognition} address a similar problem by combining monocular camera recognition with ITS information. Their approach demonstrates the value of using external traffic infrastructure information to support camera-based recognition. In this project, the same principle is adapted using GNSS/GPS and a traffic-light node database. The GPS module estimates the current vehicle position, while the node database provides the locations of signalized intersections. By calculating the distance between the vehicle and the nearest traffic-light node, the system can determine whether the vehicle is inside the warning range and whether the AI perception pipeline should be activated.

Compared with a camera-only warning system, this hybrid design reduces unnecessary inference when the vehicle is far from relevant intersections and improves the chance of producing an early warning when the vehicle approaches a traffic light. Thus, spatial positioning is not used as a replacement for visual perception, but as a triggering and contextual support mechanism that makes the perception pipeline more practical for embedded ADAS deployment.

Based on the gaps identified in these related works, the next section consolidates the chapter and prepares the transition toward the theoretical background required for the proposed implementation.

